%--- E14: Section 7.3, after the dictionary criterion ------------------------
\begin{explained}{identifying a nonlinear equation relative to a dictionary}
\textbf{Step 1: why $\Ann(u)$ is no longer enough.} For a linear equation, e.g.\ $u=e^{x-t}$, $(\partial_t+\partial_x)u=0$ is the action of a linear constant-coefficient operator. In $u_t+uu_x=0$ the term $uu_x$ multiplies two quantities depending on $u$: it is not the action of a linear constant-coefficient operator on $u$, so we need another language: polynomial expressions in the jet variables.

\medskip\textbf{Step 2: jets and the dictionary.} Fix a maximal order $k$ and take the jet variables $u,u_t,u_x,u_{tt},u_{tx},u_{xx},\dots$ up to order $k$. A polynomial in them is a candidate differential relation; Burgers is $F(u,u_t,u_x)=u_t+uu_x$. A \emph{dictionary} $\mathcal D$ is a finite list of candidate monomials, e.g.\ $\{u_t,u_x,u^2,uu_x,u_{xx},\dots\}$, and $V=\spn\mathcal D$ contains all linear combinations. A candidate equation is $c_1u_t+c_2uu_x+c_3u_{xx}=0$; identifying the equation from data means finding the coefficients for which the relation holds on the observed solution. $V$ is a finite-dimensional vector space of candidates, \emph{not} necessarily an ideal: the dictionary limits what one is willing to consider.

\medskip\textbf{Step 3: $I(u)$ versus $[F]$.} $I(u)$ is the differential ideal of \emph{all} polynomial differential relations satisfied by $u$. If $F=u_t+uu_x=0$ holds identically, so do its derivatives, e.g.\ $D_xF=u_{tx}+u_x^2+uu_{xx}$ (total derivative $D_x$). $[F]$, the differential ideal generated by $F$, contains $F$, its differential consequences $D_xF,D_tF,\dots$ and polynomial combinations of them. If $u$ solves $F=0$ then $[F]\subseteq I(u)$, but the inclusion can be strict: a particular solution may satisfy additional relations that do not follow from the equation alone. That is the core of identifiability.

\medskip\textbf{Step 4: why ``$I(u)\cap V=\langle F\rangle$'' is wrong.} If $\mathcal D$ contains the monomials of both $F$ and $D_xF$, both lie in $V$. If $\langle F\rangle$ meant only scalar multiples of $F$, it would not contain $D_xF$, which is certainly a relation satisfied by the solution. The correct comparison is
\[
I(u)\cap V=[F]\cap V:
\]
every relation of the dictionary that the data satisfy is already a differential consequence of the proposed equation. It does \emph{not} require the data to satisfy only the original equation and nothing derived from it.

\medskip\textbf{Step 5: the exact sequence.}
\[
0\to[F]\cap V\to I(u)\cap V\to\mathcal J_u\cap\psi(V)\to0 .
\]
$\mathcal J_u$ is the jet-level obstruction (distinct from the cofactor $J_u$ of Section~\ref{sec:prelim}) and $\psi$ the map from candidate expressions to jet-level objects (its precise definition is flagged in the author's notes). The first term holds the dictionary relations explained by the proposed equation, the middle term every dictionary relation the data satisfy, the last the extra relations left after quotienting out the consequences of $F$. Exactness gives $I(u)\cap V=[F]\cap V\iff\mathcal J_u\cap\psi(V)=\{0\}$: the dictionary contains no detectable obstruction. It is a \emph{finite-dictionary} criterion: it does not say the data identify the equation among all conceivable nonlinear PDEs, only relative to the chosen candidate space.

\medskip\textbf{Step 6: shocks and solitons.} Identifiability depends on the solution, not only on the equation. \emph{Burgers with $N$ shocks:} in the paper's evolutionary dictionary the equation is identified iff the maximal polynomial degree satisfies $d\le N$; shock structure gives information that distinguishes the equation from competing polynomial relations, depending on the number of shocks and on the allowed degree of candidate terms (specific to the dictionary and assumptions of the cited result). \emph{$N$-soliton KdV:} the identifying polynomial has total degree $(N+2)2^{N-1}$; for $N=2$ this is $(2+2)\cdot2^{1}=8$. A special solution can satisfy many polynomial relations, so a sufficiently rich dictionary is needed to reveal which relation actually identifies the equation rather than describing that particular solution.

\medskip\textbf{Step 7: why a closed form can fail to identify.} $u=e^{x+t}$ satisfies $u_t-u=0$, $u_x-u=0$, and consequences such as $u_{tx}-u=0$, $u_{xx}-u=0$. The solution alone does not say which relation is ``the'' governing equation: several equations explain the same special function. That is the intuition behind ``closed form $\Rightarrow$ does not identify''; the exact meaning of ``closed form'' in the theorem is the one fixed in the cited result, but the broad lesson is that an explicit solution may be too special to separate the true equation from competitors.

\medskip\textbf{The big picture.} The question changes from ``which linear differential operators annihilate $u$?'' to ``which polynomial differential relations does $u$ satisfy, and which genuinely identify its equation rather than merely describing a special solution?''. The dictionary makes the question finite-dimensional; $[F]$ accounts for all consequences of the proposed equation; $\mathcal J_u\cap\psi(V)$ detects extra relations. The crucial distinction is between \emph{fitting} a solution and \emph{identifying} its equation.
\end{explained}

